%%%%%%%%%%%%%%%%%%%%%%%%%%%%%%%%%%%%%%%%%%%%%%%%%%%%%%%%%%%%%%%%%%%%%%%%%%%%%%%
% APPENDIX AG: Complexity Economics
% Emergence, Adaptation, and Non-Equilibrium Dynamics
% Category: DOMAIN | Primary Chapter: 12 | Last Updated: 2026-01-11
% Language: English
%%%%%%%%%%%%%%%%%%%%%%%%%%%%%%%%%%%%%%%%%%%%%%%%%%%%%%%%%%%%%%%%%%%%%%%%%%%%%%%
% -----------------------------------------------------------------------------
% APPENDIX SCOPE BOX
% -----------------------------------------------------------------------------

\begin{tcolorbox}[
    colback=orange!5!white,
    colframe=orange!75!black,
    title={\textbf{Appendix Scope: Ziel / In-Scope / Out-of-Scope / Constraints}},
    fonttitle=\bfseries
]

\textbf{Ziel (Objective):}
\begin{itemize}[nosep]
    \item [TODO: Define objective for Unknown Title]
\end{itemize}

\textbf{In-Scope:}
\begin{itemize}[nosep]
    \item Complexity Economics: Beyond Equilibrium
    \item Worked Example
    \item Results
\end{itemize}

\textbf{Out-of-Scope (delegiert an andere Appendices):}
\begin{itemize}[nosep]
    \item Glossary $\rightarrow$ Appendix GLS
    \item Behavioral GT $\rightarrow$ Appendix SPC
    \item CORE-HOW $\rightarrow$ Appendix HOW
    \item Behavioral GT $\rightarrow$ Appendix SPC
    \item Awareness $\rightarrow$ Appendix AWA
\end{itemize}

\textbf{Constraints (Anwendungsgrenzen):}
\begin{itemize}[nosep]
    \item [TODO: Define when this appendix does NOT apply]
    \item [TODO: Define required prerequisites]
    \item [TODO: Define known limitations]
\end{itemize}

\textbf{Lieferobjekte (Deliverables):}
\begin{itemize}[nosep]
    \item 42 Equation(s)
    \item Worked Example(s)
\end{itemize}

\end{tcolorbox}


\section{Complexity Economics: Beyond Equilibrium}
\label{app:complexity-economics}

%==============================================================================
% FORMAL COMPLIANCE BLOCK
%==============================================================================

% --- Header Block ---
\begin{tcolorbox}[colback=gray!5!white, colframe=gray!75!black,
    title=\textbf{Appendix: AG DOMAIN-COMPLEXITY}]
\textbf{Category:} DOMAIN (Application Field) \\
\textbf{Title:} Complexity Economics: Beyond Equilibrium \\
\textbf{Core Question:} How do equilibria emerge, evolve, and sometimes collapse? \\
\textbf{Primary Chapter:} Ch.\ 12 (Integration) \\
\textbf{Last Updated:} 2026-01-11
\end{tcolorbox}

% --- Abstract ---
\begin{tcolorbox}[colback=blue!5!white, colframe=blue!75!black,
    title=\textbf{Abstract}]

This appendix synthesizes \textbf{complexity economics}---the study of economies
as complex adaptive systems where equilibria emerge from agent interactions,
are path-dependent, and perpetually evolving. Key insights:
(1) \textbf{Emergence}: Macro patterns arise from micro interactions without central design;
(2) \textbf{Path dependence}: History matters---early events shape long-run outcomes;
(3) \textbf{Heterogeneity}: Agents differ, representative agent fails (Kirman 1992);
(4) \textbf{Non-equilibrium}: Economy is perpetually adapting, not at rest.

The Beauty Contest generative framework (M\&N 2018) shows that heterogeneous
Level-$k$ reasoning is \textit{structured heterogeneity}---not noise---providing
cognitive microfoundations for complexity dynamics.

\textbf{Core Contribution:} 20 foundational papers with 80,000+ combined citations
establishing complexity as essential lens for understanding $C^*$ emergence.

\end{tcolorbox}

% --- Terminology Reference ---
\begin{tcolorbox}[colback=gray!5!white, colframe=gray!50!black]
\textbf{Terminology:} See Appendix GLS (Glossary) for standard definitions.
Key terms: emergence, path dependence, increasing returns, phase transition,
self-organized criticality, heterogeneous agents, complex adaptive system.
\end{tcolorbox}

% --- Chapter Linkage Box ---
\begin{tcolorbox}[colback=green!5!white, colframe=green!75!black,
    title=\textbf{Chapter References}]

\textbf{Primary Chapter:} Ch.\ 12 (Integration)
\begin{itemize}[nosep]
    \item Complexity as integrating framework for behavioral economics
    \item Emergence of coherent behavior ($K$)
\end{itemize}

\textbf{Secondary Chapters:}
\begin{itemize}[nosep]
    \item Ch.\ 5 (Complementarity): $\gamma > 0$ creates increasing returns
    \item Ch.\ 9 (Context): $\Psi$ as endogenous, path-dependent
    \item Ch.\ 11 (Predictions): Non-linear dynamics, phase transitions
\end{itemize}

\textbf{Cross-References:}
\begin{itemize}[nosep]
    \item Appendix SPC (Behavioral GT): Level-$k$ heterogeneity
    \item Appendix MSC Section 11: M\&N structured heterogeneity
    \item Appendix HOW (CORE-HOW): Complementarity creates feedback loops
\end{itemize}

\end{tcolorbox}

% --- Quick Reference ---
\begin{tcolorbox}[colback=yellow!5!white, colframe=yellow!75!black,
    title=\textbf{Quick Reference: Key Concepts}]

\renewcommand{\arraystretch}{1.3}
\begin{tabular}{@{}p{4cm}p{5cm}p{4cm}@{}}
\toprule
\textbf{Concept} & \textbf{Formulation} & \textbf{EBF Element} \\
\midrule
Emergence & $C^*_{macro} \neq \sum C^*_{micro}$ & Aggregation fails \\
Path dependence & $C^*_t = f(C^*_{<t})$ & History matters \\
Increasing returns & $\gamma > 0 \Rightarrow$ Lock-in & Complementarity \\
Phase transition & Small $\Delta\Psi \to$ Large $\Delta C^*$ & Non-linearity \\
Heterogeneous agents & $\Psi_i \neq \Psi_j$ & No representative agent \\
\bottomrule
\end{tabular}

\end{tcolorbox}

% --- Bibliography Reference ---
\noindent\textbf{Full citations:} See \texttt{bibliography/bcm\_master.bib} (Master Bibliography) \\
\textbf{Primary sources:} Arthur (1994), Schelling (1978), Kirman (1992)

\vspace{1em}

%==============================================================================
% PART 0: INTEGRATION OVERVIEW
%==============================================================================
\subsection{Framework Integration Map}

This appendix demonstrates how complexity economics provides essential 
foundations for the $(C, \Psi, C^*, K, Q)$ framework. Standard economics 
assumes equilibrium exists, is unique, and is reached quickly. Complexity 
economics shows that equilibria emerge from agent interactions, may be 
multiple, path-dependent, and perpetually adapting. This fundamentally 
changes how we understand $C^*$: not as a fixed point to be calculated, 
but as an emergent property of complex adaptive systems.

\begin{center}
\renewcommand{\arraystretch}{1.4}
\begin{tabular}{|l|l|l|}
\hline
\textbf{Framework Element} & \textbf{Complexity Contribution} & \textbf{Key Papers} \\
\hline
$C^*$ as emergent & Not computed, emerges & Arthur 1994, 2015 \\
Path dependence & $C^*_t = f(C^*_{<t})$ & Arthur 1989, David 1985 \\
$C^*_{macro} \neq \sum C^*_{micro}$ & Emergence & Schelling 1971, 1978 \\
Heterogeneous $\Psi_i$ & No representative agent & Kirman 1992 \\
Adaptive dynamics & Learning, evolution & Holland 1992, Epstein 1996 \\
Non-equilibrium & Perpetual novelty & Farmer 2009, Bak 1996 \\
\hline
\end{tabular}
\end{center}

\paragraph{Complexity Economics' Unique Role.}
Complexity economics contributes uniquely to the framework:
\begin{enumerate}
    \item \textbf{Emergence}: Macro patterns from micro interactions
    \item \textbf{Path dependence}: History matters for $C^*$
    \item \textbf{Heterogeneity}: Agents differ, aggregation fails
    \item \textbf{Adaptation}: Agents learn, system evolves
    \item \textbf{Non-equilibrium}: Economy perpetually in flux
    \item \textbf{Non-linearity}: Small causes, large effects
\end{enumerate}

%==============================================================================
% PART I: FOUNDATIONS – COMPLEXITY SCIENCE
%==============================================================================
\subsection{Part I: Foundations – The Complexity Paradigm}

\subsubsection{What is Complexity Economics?}

Arthur (2015) defines complexity economics:

\begin{equation}
\boxed{
\text{Economy} = \text{Complex Adaptive System of heterogeneous interacting agents}
}
\label{eq:complexity-def}
\end{equation}

\paragraph{Key Features.}
\begin{itemize}
    \item \textbf{Heterogeneous agents}: Not identical, not representative
    \item \textbf{Local interaction}: Agents interact with neighbors, not all
    \item \textbf{Bounded rationality}: $\Psi_C < \infty$
    \item \textbf{Adaptation}: Agents learn, strategies evolve
    \item \textbf{Emergence}: Macro patterns not in micro rules
    \item \textbf{Non-equilibrium}: Perpetual change, novelty
\end{itemize}

\subsubsection{Santa Fe Institute}

The Santa Fe Institute (founded 1984) brought complexity science to economics:

\begin{center}
\renewcommand{\arraystretch}{1.3}
\begin{tabular}{|l|l|}
\hline
\textbf{Physicist/Scientist} & \textbf{Economic Contribution} \\
\hline
Kenneth Arrow & Founding vision, general equilibrium limits \\
Philip Anderson & ``More is Different'' --- emergence \\
W. Brian Arthur & Increasing returns, path dependence \\
John Holland & Genetic algorithms, adaptation \\
Stuart Kauffman & Self-organization, NK landscapes \\
\hline
\end{tabular}
\end{center}

\subsubsection{Simon: The Architecture of Complexity}

Herbert Simon (1962) laid foundations:

\begin{equation}
\boxed{
\text{Complex systems are nearly decomposable hierarchies}
}
\label{eq:simon-hierarchy}
\end{equation}

\paragraph{Key Insights.}
\begin{itemize}
    \item Hierarchical structure enables evolution
    \item Bounded rationality is necessary, not pathological
    \item Satisficing replaces maximizing
\end{itemize}

\subsubsection{Framework Translation}

Simon shows $\Psi_C$ limits are fundamental:
\begin{equation}
\Psi_C^{bounded} \Rightarrow \text{Satisficing, not optimizing}
\end{equation}

%==============================================================================
% PART II: INCREASING RETURNS AND PATH DEPENDENCE
%==============================================================================
\subsection{Part II: Increasing Returns and Path Dependence}

\subsubsection{Arthur's Breakthrough}

W. Brian Arthur challenged the core assumption:

\begin{equation}
\boxed{
\text{Increasing returns} \Rightarrow \text{Multiple equilibria + Path dependence + Lock-in}
}
\label{eq:arthur-ir}
\end{equation}

\paragraph{Standard Economics.}
Diminishing returns $\Rightarrow$ unique equilibrium, history doesn't matter.

\paragraph{Arthur's Alternative.}
Increasing returns $\Rightarrow$ many equilibria, history selects which one.

\subsubsection{Sources of Increasing Returns}

Arthur (1994) identifies sources:
\begin{enumerate}
    \item \textbf{Learning by doing}: Costs fall with cumulative production
    \item \textbf{Network effects}: Value rises with users ($C_{users}^{network} > 0$)
    \item \textbf{Economies of scale}: Fixed costs spread over output
    \item \textbf{Informational returns}: Knowledge spillovers
    \item \textbf{Technological interrelatedness}: Complementary investments
\end{enumerate}

\subsubsection{Path Dependence}

With increasing returns:
\begin{equation}
\boxed{
C^*_T = f(C_0, C_1, \ldots, C_{CAM-1}, \text{random events})
}
\label{eq:path-dependence}
\end{equation}

Small early events can have large long-run consequences.

\paragraph{QWERTY Example.}
Keyboard layout persists despite inefficiency due to:
\begin{itemize}
    \item Early adoption by typing schools
    \item Network effects (typists trained on QWERTY)
    \item Switching costs
\end{itemize}

\subsubsection{Framework Translation}

Path dependence means $C^*$ is historically contingent:
\begin{equation}
C^*_{selected} \in \{C^*_1, C^*_2, \ldots\} \quad \text{selected by history}
\end{equation}

Multiple equilibria, path selects one.

\subsubsection{Lock-In and Inefficiency}

Path dependence can lead to inferior outcomes:
\begin{equation}
C^*_{selected} \neq C^*_{efficient}
\end{equation}

Early random events may lock in suboptimal technology/institution.

%==============================================================================
% PART III: EMERGENCE AND SELF-ORGANIZATION
%==============================================================================
\subsection{Part III: Emergence and Self-Organization}

\subsubsection{Schelling's Segregation Model}

Schelling (1971, 1978) demonstrated emergence:

\begin{equation}
\boxed{
\text{Mild individual preferences} \Rightarrow \text{Complete segregation}
}
\label{eq:schelling}
\end{equation}

\paragraph{The Model.}
\begin{itemize}
    \item Agents on grid, two types
    \item Each prefers $\geq 1/3$ neighbors same type (mild preference)
    \item Unhappy agents move to random vacant spot
\end{itemize}

\paragraph{Result.}
Complete segregation emerges, far exceeding individual preferences!

\subsubsection{Framework Translation}

Emergence means macro $\neq$ sum of micro:
\begin{equation}
\boxed{
C^*_{macro} \neq \sum_i C^*_i \quad \text{(emergent properties)}
}
\label{eq:emergence}
\end{equation}

Cannot predict macro from micro without simulation.

\subsubsection{Tipping Points}

Schelling identified tipping dynamics:
\begin{equation}
\text{Small change in } \Psi \Rightarrow \text{Large discontinuous change in } C^*
\end{equation}

\paragraph{Examples.}
\begin{itemize}
    \item Neighborhood tipping (white flight)
    \item Bank runs (confidence collapse)
    \item Technology adoption (network effects)
    \item Social norms (critical mass)
\end{itemize}

\subsubsection{Self-Organized Criticality}

Bak (1996) showed many systems self-organize to criticality:

\begin{equation}
\boxed{
\text{System} \to \text{Critical state} \Rightarrow \text{Power-law distributed events}
}
\label{eq:soc}
\end{equation}

\paragraph{Sandpile Model.}
Drop sand grains; avalanches of all sizes occur with power-law frequency.

\paragraph{Economic Implications.}
\begin{itemize}
    \item Large crashes are not ``outliers'' but inherent
    \item Cannot predict individual events
    \item System statistics are predictable
\end{itemize}

%==============================================================================
% PART IV: AGENT-BASED MODELING
%==============================================================================
\subsection{Part IV: Agent-Based Modeling – Growing $C^*$}

\subsubsection{The ABM Approach}

Agent-based modeling (ABM) simulates emergent $C^*$:

\begin{equation}
\boxed{
\text{Define agents + rules} \to \text{Simulate} \to \text{Observe emergent } C^*
}
\label{eq:abm}
\end{equation}

\paragraph{Epstein's Dictum.}
``If you didn't grow it, you didn't explain it.''

\subsubsection{Holland: Genetic Algorithms and Adaptation}

John Holland (1992) introduced adaptive agents:

\begin{equation}
\text{Agent strategies evolve via selection, crossover, mutation}
\end{equation}

\paragraph{Classifier Systems.}
\begin{itemize}
    \item Rules compete based on past success
    \item Successful rules reproduce
    \item New rules created by recombination
\end{itemize}

\subsubsection{Framework Translation}

ABM shows how $C^*$ emerges:
\begin{equation}
C^*_{ABM} = \text{Emergent pattern from agent interactions}
\end{equation}

Not solved analytically, but ``grown'' computationally.

\subsubsection{Sugarscape: Growing Artificial Societies}

Epstein \& Axtell (1996) created Sugarscape:

\begin{equation}
\text{Simple agents} + \text{Simple rules} \Rightarrow \text{Complex social phenomena}
\end{equation}

\paragraph{Emergent Phenomena.}
\begin{itemize}
    \item Trade and markets
    \item Wealth inequality (Pareto distribution!)
    \item Migration patterns
    \item Cultural diffusion
    \item Combat and disease spread
\end{itemize}

\subsubsection{Tesfatsion: ACE Methodology}

Tesfatsion (2006) systematized Agent-based Computational Economics:

\begin{equation}
\text{ACE} = \text{Computational laboratory for testing economic theories}
\end{equation}

\paragraph{Advantages.}
\begin{itemize}
    \item Heterogeneous agents
    \item Bounded rationality
    \item Local interaction
    \item Out-of-equilibrium dynamics
    \item Emergent macro patterns
\end{itemize}

%==============================================================================
% PART V: HETEROGENEITY AND EXPECTATIONS
%==============================================================================
\subsection{Part V: Heterogeneity – The Death of Representative Agent}

\subsubsection{Kirman's Critique}

Kirman (1992) devastatingly critiqued representative agent:

\begin{equation}
\boxed{
\text{Representative agent models have no theoretical foundation}
}
\label{eq:kirman}
\end{equation}

\paragraph{Problems.}
\begin{itemize}
    \item Aggregate behavior $\neq$ scaled individual behavior
    \item Sonnenschein-Mantel-Debreu: Any aggregate demand possible
    \item Heterogeneity creates dynamics absent from representative agent
\end{itemize}

\subsubsection{The Ants Analogy}

Kirman's ants model:
\begin{itemize}
    \item Two food sources, identical
    \item Ants follow pheromone trails
    \item Aggregate behavior: Wild swings between sources
    \item Individual behavior: Simple following rule
\end{itemize}

\paragraph{Lesson.}
Aggregate volatility emerges from interaction, not from individual volatility.

\subsubsection{Framework Translation}

Heterogeneity is fundamental to $\Psi$:
\begin{equation}
\Psi = \{\Psi_1, \Psi_2, \ldots, \Psi_n\} \quad \text{not } \Psi_{representative}
\end{equation}

Distribution of $\Psi$ matters, not just mean.

\subsubsection{Hommes: Heterogeneous Expectations}

Hommes (2006) models heterogeneous expectations:

\begin{equation}
\text{Agents use different forecasting rules, switch based on performance}
\end{equation}

\paragraph{Results.}
\begin{itemize}
    \item Fundamentalists vs. chartists
    \item Endogenous bubbles and crashes
    \item Excess volatility without exogenous shocks
    \item Fat tails emerge from interaction
\end{itemize}

\subsubsection{Framework Translation}

Heterogeneous $\Psi_K$ (expectations) creates dynamics:
\begin{equation}
\Psi_K^{heterogeneous} \Rightarrow \text{Endogenous volatility, bubbles}
\end{equation}

\subsubsection{Mauersberger \& Nagel: Structured Heterogeneity}

Mauersberger \& Nagel (2018) provide cognitive microfoundations for heterogeneity:

\begin{tcolorbox}[colback=yellow!5!white, colframe=yellow!75!black,
    title=\textbf{M\&N (2018): Level-$k$ as Structured Heterogeneity}]

\textbf{Key Insight for Complexity Economics:}
\begin{quote}
``Heterogeneity is \textbf{structured}, not noise.''
\end{quote}

\textbf{Level-$k$ Distribution:}
\begin{itemize}[nosep]
    \item Level-0: $\sim$20\% (naive, anchor-driven)
    \item Level-1: $\sim$40\% (best-respond to naive)
    \item Level-2: $\sim$30\% (second-order reasoning)
    \item Level-3+: $\sim$10\% (approaching equilibrium)
\end{itemize}

This is \textbf{not noise}---it is systematic cognitive heterogeneity that:
\begin{enumerate}[nosep]
    \item Provides microfoundation for Kirman's critique
    \item Explains Hommes' fundamentalist/chartist distribution
    \item Generates endogenous dynamics without representative agent
\end{enumerate}

\end{tcolorbox}

\paragraph{Framework Translation.}
Level-$k$ connects to EBF's Awareness $A(\cdot)$:
\begin{equation}
\Psi_C^{Level-k} \leftrightarrow A(\cdot) \quad \Rightarrow \quad
\text{Population distribution of } A \text{ drives complexity}
\end{equation}

\textbf{Cross-Reference:} Appendix MSC Section 11, Appendix SPC (Behavioral GT),
Appendix AWA (Awareness).

%==============================================================================
% PART VI: FINANCIAL COMPLEXITY
%==============================================================================
\subsection{Part VI: Financial Complexity – Markets as Complex Systems}

\subsubsection{Farmer: Agent-Based Finance}

Doyne Farmer advocated ABM for finance:

\begin{equation}
\boxed{
\text{Financial markets} = \text{Ecology of trading strategies}
}
\label{eq:farmer}
\end{equation}

\paragraph{Key Insights.}
\begin{itemize}
    \item Markets never reach equilibrium
    \item Arbitrage opportunities persist
    \item Strategies interact, co-evolve
    \item Crashes are endogenous
\end{itemize}

\subsubsection{Farmer \& Foley: The Case for ABM}

Farmer \& Foley (2009) argued post-crisis:

\begin{equation}
\text{DSGE failed} \Rightarrow \text{Need ABM for macroeconomics}
\end{equation}

\paragraph{DSGE Failures.}
\begin{itemize}
    \item Missed 2008 crisis entirely
    \item Representative agent ignores contagion
    \item Equilibrium assumption ruled out crashes
    \item No financial sector in many models
\end{itemize}

\subsubsection{LeBaron: Artificial Stock Markets}

LeBaron (2006) created artificial markets:

\begin{equation}
\text{Heterogeneous traders} \to \text{Realistic market statistics}
\end{equation}

\paragraph{Emergent Properties.}
\begin{itemize}
    \item Fat tails in returns
    \item Volatility clustering
    \item Long memory
    \item Excess trading volume
\end{itemize}

All emerge without being programmed!

\subsubsection{Sornette: Dragon Kings and Crashes}

Sornette (2003) analyzed extreme events:

\begin{equation}
\boxed{
\text{Some crashes are ``Dragon Kings''} = \text{Predictable outliers}
}
\label{eq:sornette}
\end{equation}

\paragraph{Log-Periodic Power Law.}
Bubbles show characteristic acceleration before crash:
\begin{equation}
p(t) \sim (t_c - t)^{-\alpha} [1 + C \cos(\omega \ln(t_c - t) + \phi)]
\end{equation}

Crashes may be partially predictable!

\subsubsection{Taleb: Black Swans and Fat Tails}

Taleb (2007) emphasized:

\begin{equation}
\boxed{
\text{Rare events dominate} \Rightarrow \text{Normal distribution is wrong}
}
\label{eq:taleb}
\end{equation}

\paragraph{Implications.}
\begin{itemize}
    \item Standard risk models vastly underestimate tail risk
    \item VaR and similar measures are dangerous
    \item Prediction is harder than we think
    \item Robustness matters more than optimization
\end{itemize}

\subsubsection{Framework Translation}

Financial $Q$ is fat-tailed:
\begin{equation}
Q \sim \text{Power law, not Normal}
\end{equation}

Extreme events are more likely than standard models predict.

%==============================================================================
% PART VII: NETWORKS AND INTERACTION
%==============================================================================
\subsection{Part VII: Networks – The Topology of $C_{ij}$}

\subsubsection{Network Structure Matters}

Complexity economics emphasizes interaction structure:

\begin{equation}
\boxed{
\text{Network topology} \Rightarrow \text{System dynamics}
}
\label{eq:networks}
\end{equation}

\paragraph{Key Network Properties.}
\begin{itemize}
    \item Degree distribution (who connects to whom)
    \item Clustering (local density)
    \item Path length (degrees of separation)
    \item Centrality (who is important)
\end{itemize}

\subsubsection{Scale-Free Networks}

Many real networks are scale-free:
\begin{equation}
P(k) \sim k^{-\gamma} \quad \text{(degree distribution)}
\end{equation}

\paragraph{Implications.}
\begin{itemize}
    \item Hubs exist (highly connected nodes)
    \item Robust to random failure
    \item Fragile to targeted attack on hubs
\end{itemize}

\subsubsection{Framework Translation}

$C_{ij}$ has network structure:
\begin{equation}
C_{ij} \neq 0 \quad \text{only for connected } (i,j)
\end{equation}

Network topology shapes complementarities.

\subsubsection{Contagion and Cascades}

Networks enable cascading failures:
\begin{equation}
\text{Local shock} \xrightarrow{\text{network}} \text{Global crisis}
\end{equation}

\paragraph{Examples.}
\begin{itemize}
    \item Financial contagion (2008)
    \item Supply chain disruptions
    \item Information cascades
    \item Epidemic spread
\end{itemize}

%==============================================================================
% PART VIII: SYNTHESIS
%==============================================================================
\subsection{Part VIII: Synthesis – Complexity Economics Equations}

\subsubsection{Complete Framework Translation}

\begin{center}
\renewcommand{\arraystretch}{1.5}
\begin{tabular}{|l|c|l|}
\hline
\textbf{Concept} & \textbf{Equation} & \textbf{Framework Element} \\
\hline
Emergence & $C^*_{macro} \neq \sum C^*_{micro}$ & Macro from micro \\
\hline
Path dependence & $C^*_T = f(C_0, \ldots, C_{CAM-1})$ & History matters \\
\hline
Increasing returns & $C_{ij} > 0$ self-reinforcing & Multiple $C^*$ \\
\hline
Tipping & Small $\Delta\Psi \to$ large $\Delta C^*$ & Discontinuity \\
\hline
Heterogeneity & $\Psi_i \neq \Psi_j$ & Distribution matters \\
\hline
Adaptation & $C_t \to C^*$ via learning & Evolutionary $C^*$ \\
\hline
Fat tails & $Q \sim$ Power law & Extreme events likely \\
\hline
Networks & $C_{ij}$ topology & Contagion, cascades \\
\hline
\end{tabular}
\end{center}

\subsubsection{The Complexity Worldview}

\begin{equation}
\boxed{
\text{Economy} = \text{Open, evolving, non-equilibrium complex adaptive system}
}
\end{equation}

Key differences from standard economics:
\begin{enumerate}
    \item \textbf{Equilibrium}: Emergent and shifting, not assumed
    \item \textbf{Agents}: Heterogeneous and adaptive, not identical and rational
    \item \textbf{Dynamics}: Perpetual change, not comparative statics
    \item \textbf{Prediction}: Pattern, not point prediction
    \item \textbf{Policy}: Adaptive, not optimal control
\end{enumerate}

\subsubsection{Complexity Economics' Contribution to Framework}

\textbf{Without Complexity Economics:}
\begin{itemize}
    \item Assume unique stable $C^*$
    \item Ignore path dependence
    \item Use representative agent
    \item Miss emergent phenomena
    \item Underestimate tail risk
\end{itemize}

\textbf{With Complexity Economics:}
\begin{itemize}
    \item $C^*$ as emergent, path-dependent
    \item Heterogeneity fundamental
    \item Adaptation and learning central
    \item Non-linearity and tipping points
    \item Networks and contagion
    \item Robustness over optimization
\end{itemize}

%==============================================================================
% PART IX: COMPLETE PAPER DOCUMENTATION
%==============================================================================
\subsection{Part IX: Complete Paper Documentation}

%--- ARTHUR ---
\subsubsection{W. Brian Arthur}

\paragraph{Paper 1: Arthur (1989)}

``Competing Technologies, Increasing Returns, and Lock-In by Historical Events.''
\textit{Economic Journal}, 99(394), 116--131.

\textbf{Summary.}
Shows how increasing returns lead to path dependence and lock-in.
Early adopters create positive feedback; market may lock into inferior technology.

\textbf{Framework Translation.}
\begin{equation}
C_{ij}^{increasing} \Rightarrow C^*_{selected} \neq C^*_{efficient} \text{ possible}
\end{equation}

\textbf{Key Finding.}
4,000+ citations. Foundational for path dependence.

\paragraph{Paper 2: Arthur (1994)}

\textit{Increasing Returns and Path Dependence in the Economy.}
University of Michigan Press.

\textbf{Summary.}
Collected essays on increasing returns. Covers technology adoption, 
positive feedbacks, self-reinforcing mechanisms.

\textbf{Framework Translation.}
Complete theory of increasing returns economics.

\paragraph{Paper 3: Arthur (2015)}

\textit{Complexity and the Economy.}
Oxford University Press.

\textbf{Summary.}
Synthesizes 25 years of complexity economics. Economy as complex 
adaptive system, perpetually evolving, not in equilibrium.

\textbf{Framework Translation.}
\begin{equation}
\text{Economy} \neq \text{Equilibrium system}
\end{equation}

\textbf{Key Finding.}
Definitive statement of complexity economics.

\paragraph{Paper 4: Arthur, Durlauf \& Lane (1997)}

\textit{The Economy as an Evolving Complex System II.}
Addison-Wesley (Santa Fe Institute).

\textbf{Summary.}
Second Santa Fe volume on complexity economics. Covers agent-based modeling, 
increasing returns, institutions, and methodology.

\textbf{Framework Translation.}
Santa Fe Institute's foundational collection.

%--- SCHELLING ---
\subsubsection{Thomas Schelling}

\paragraph{Paper 5: Schelling (1971)}

``Dynamic Models of Segregation.''
\textit{Journal of Mathematical Sociology}, 1(2), 143--186.

\textbf{Summary.}
Shows complete segregation emerges from mild individual preferences.
Classic demonstration of emergence.

\textbf{Framework Translation.}
\begin{equation}
\text{Mild } \Psi_i \Rightarrow \text{Extreme } C^*_{macro}
\end{equation}

\textbf{Key Finding.}
Nobel Prize (2005). 5,000+ citations.

\paragraph{Paper 6: Schelling (1978)}

\textit{Micromotives and Macrobehavior.}
W.W. Norton.

\textbf{Summary.}
Explores how individual behavior aggregates to unexpected macro patterns.
Tipping, critical mass, self-fulfilling prophecies.

\textbf{Framework Translation.}
Micro-macro gap is fundamental.

%--- SIMON ---
\subsubsection{Herbert Simon}

\paragraph{Paper 7: Simon (1962)}

``The Architecture of Complexity.''
\textit{Proceedings of the American Philosophical Society}, 106(6), 467--482.

\textbf{Summary.}
Complex systems are hierarchical and nearly decomposable.
Evolution builds complexity through stable subassemblies.

\textbf{Framework Translation.}
$C$ has hierarchical structure enabling evolution.

\textbf{Key Finding.}
Nobel Prize (1978). Foundation of complexity science.

%--- HOLLAND AND KAUFFMAN ---
\subsubsection{Adaptation and Self-Organization}

\paragraph{Paper 8: Holland (1992)}

\textit{Adaptation in Natural and Artificial Systems.}
MIT Press. (Originally 1975)

\textbf{Summary.}
Introduces genetic algorithms. Shows how adaptation works through 
selection, crossover, and mutation.

\textbf{Framework Translation.}
$C^*$ evolves through adaptive selection.

\paragraph{Paper 9: Holland (1995)}

\textit{Hidden Order: How Adaptation Builds Complexity.}
Addison-Wesley.

\textbf{Summary.}
How complex adaptive systems work. Agents, tags, internal models, 
building blocks.

\textbf{Framework Translation.}
Complexity from simple adaptive rules.

\paragraph{Paper 10: Kauffman (1993)}

\textit{The Origins of Order: Self-Organization and Selection in Evolution.}
Oxford University Press.

\textbf{Summary.}
Self-organization complements natural selection. NK landscapes, 
edge of chaos, spontaneous order.

\textbf{Framework Translation.}
$C^*$ landscapes have structure affecting evolution.

\paragraph{Paper 11: Bak (1996)}

\textit{How Nature Works: The Science of Self-Organized Criticality.}
Copernicus.

\textbf{Summary.}
Many systems self-organize to critical state. Power-law distributions, 
sandpile model, earthquakes, evolution.

\textbf{Framework Translation.}
\begin{equation}
\text{System} \to \text{Critical} \Rightarrow P(\text{event size } s) \sim s^{-\alpha}
\end{equation}

%--- HETEROGENEITY ---
\subsubsection{Heterogeneity and Expectations}

\paragraph{Paper 12: Kirman (1992)}

``Whom or What Does the Representative Individual Represent?''
\textit{Journal of Economic Perspectives}, 6(2), 117--136.

\textbf{Summary.}
Devastating critique of representative agent. No theoretical foundation; 
aggregate behavior differs qualitatively from individual.

\textbf{Framework Translation.}
\begin{equation}
\nexists \Psi_{representative}: C^*(\Psi_{rep}) = C^*(\{\Psi_i\})
\end{equation}

\textbf{Key Finding.}
2,500+ citations. Essential methodological critique.

\paragraph{Paper 13: Hommes (2006)}

``Heterogeneous Agent Models in Economics and Finance.''
\textit{Handbook of Computational Economics}, Vol. 2.

\textbf{Summary.}
Survey of heterogeneous agent models. Expectations switching, 
endogenous dynamics, bubbles, crashes.

\textbf{Framework Translation.}
Heterogeneous $\Psi_K$ creates endogenous dynamics.

%--- ABM ---
\subsubsection{Agent-Based Modeling}

\paragraph{Paper 14: Epstein \& Axtell (1996)}

\textit{Growing Artificial Societies: Social Science from the Bottom Up.}
MIT Press.

\textbf{Summary.}
Sugarscape model. Grows trade, warfare, culture, disease from 
simple agents on simple landscape.

\textbf{Framework Translation.}
``If you didn't grow it, you didn't explain it.''

\textbf{Key Finding.}
5,000+ citations. ABM classic.

\paragraph{Paper 15: Epstein (2006)}

\textit{Generative Social Science: Studies in Agent-Based Computational Modeling.}
Princeton University Press.

\textbf{Summary.}
Methodology of generative explanation. Applications to civil violence, 
epidemics, retirement, norms.

\textbf{Framework Translation.}
Generative sufficiency as explanatory standard.

\paragraph{Paper 16: Tesfatsion (2006)}

``Agent-Based Computational Economics: A Constructive Approach to Economic Theory.''
\textit{Handbook of Computational Economics}, Vol. 2.

\textbf{Summary.}
Comprehensive survey of ACE methodology. Design principles, 
applications, validation.

\textbf{Framework Translation.}
ACE as computational laboratory for $C^*$ emergence.

%--- FINANCIAL COMPLEXITY ---
\subsubsection{Financial Complexity}

\paragraph{Paper 17: Farmer \& Foley (2009)}

``The Economy Needs Agent-Based Modelling.''
\textit{Nature}, 460, 685--686.

\textbf{Summary.}
Post-crisis call for ABM in macroeconomics. DSGE failed; 
need models with heterogeneity and non-equilibrium.

\textbf{Framework Translation.}
Equilibrium models insufficient for crisis.

\paragraph{Paper 18: LeBaron (2006)}

``Agent-Based Computational Finance.''
\textit{Handbook of Computational Economics}, Vol. 2.

\textbf{Summary.}
Survey of artificial stock markets. Heterogeneous traders generate 
realistic market statistics endogenously.

\textbf{Framework Translation.}
Fat tails, volatility clustering emerge from interaction.

\paragraph{Paper 19: Sornette (2003)}

\textit{Why Stock Markets Crash: Critical Events in Complex Financial Systems.}
Princeton University Press.

\textbf{Summary.}
Crashes as critical phenomena. Log-periodic power laws, dragon kings, 
partial predictability of bubbles.

\textbf{Framework Translation.}
\begin{equation}
\text{Bubbles show characteristic patterns before crash}
\end{equation}

\paragraph{Paper 20: Taleb (2007)}

\textit{The Black Swan: The Impact of the Highly Improbable.}
Random House.

\textbf{Summary.}
Rare events dominate history. Normal distribution vastly underestimates 
tail risk. Need robustness, not optimization.

\textbf{Framework Translation.}
\begin{equation}
Q \sim \text{Fat tails} \Rightarrow \text{Expect the unexpected}
\end{equation}

\textbf{Key Finding.}
Bestseller. Major influence on risk thinking.

%%%%%%%%%%%%%%%%%%%%%%%%%%%%%%%%%%%%%%%%%%%%%%%%%%%%%%%%%%%%%%%%%%%%%%%%%%%%%%%
% GLOBAL ORDER AS COMPLEX SYSTEM
%%%%%%%%%%%%%%%%%%%%%%%%%%%%%%%%%%%%%%%%%%%%%%%%%%%%%%%%%%%%%%%%%%%%%%%%%%%%%%%
\subsection{The Global Order as a Complex Adaptive System}
\label{sec:global-complex}

The international economic and political order exhibits all the characteristics
of a complex adaptive system: multiple interacting agents (nations, institutions,
firms), feedback loops, path dependence, and the potential for phase transitions.

\subsubsection{International Complementarity}

The global order involves complementarities at the highest aggregation level:
\begin{itemize}
    \item \textbf{Trade and peace:} Economic interdependence raises the cost of war
    \item \textbf{Institutions and cooperation:} WTO, IMF, UN enable coordination
    \item \textbf{Norms and behavior:} Shared expectations stabilize interactions
    \item \textbf{Hegemonic stability:} A dominant power provides public goods
\end{itemize}

\paragraph{Framework Translation.}
\begin{equation}
\gamma_{trade,peace} > 0, \quad \gamma_{institutions,cooperation} > 0, \quad
\gamma_{hegemony,stability} > 0
\end{equation}

\subsubsection{The Coherent Global Order: Bretton Woods}

The post-WWII Bretton Woods system exemplified global coherence ($K_{global} \approx 1$):

\begin{center}
\renewcommand{\arraystretch}{1.3}
\begin{tabular}{|l|l|}
\hline
\textbf{Element} & \textbf{Complementary Effect} \\
\hline
US hegemony & Provided public goods (security, reserve currency) \\
Fixed exchange rates & Reduced uncertainty, enabled trade \\
Capital controls & Maintained policy autonomy \\
GATT/IMF/World Bank & Institutional coordination \\
Western alliance & Security and economic alignment \\
Embedded liberalism & Domestic welfare state + international openness \\
\hline
\end{tabular}
\end{center}

This configuration was \textbf{coherent}---each element reinforced the others.
US hegemony made institutions credible; institutions enabled trade; trade
generated prosperity; prosperity supported US leadership.

\subsubsection{Current Global Incoherence}

The contemporary global order exhibits increasing incoherence ($K_{global} < 1$):

\begin{center}
\renewcommand{\arraystretch}{1.3}
\begin{tabular}{|l|l|}
\hline
\textbf{Transition} & \textbf{Manifestation} \\
\hline
US hegemony $\to$ multipolar & China rise, great power competition \\
WTO consensus $\to$ fragmentation & Trade disputes, regional blocs \\
Climate agreement $\to$ conflict & Paris vs. energy security \\
WHO authority $\to$ distrust & Pandemic response failures \\
Free trade norm $\to$ industrial policy & Reshoring, friend-shoring \\
Globalization $\to$ deglobalization & Supply chain reconfiguration \\
\hline
\end{tabular}
\end{center}

\paragraph{Complexity Interpretation.}
The global system is undergoing a \textbf{phase transition}---moving from one
coherent configuration (US hegemony) toward another (multipolar order), with
the transition period characterized by:
\begin{itemize}
    \item Increased volatility and uncertainty
    \item Norm conflict (sovereignty vs. intervention, free trade vs. industrial policy)
    \item Institutional erosion and competing alternatives
    \item Multiple possible future equilibria
\end{itemize}

\paragraph{Framework Translation.}
\begin{equation}
K_{global}^{Bretton Woods} \approx 1 \quad \to \quad K_{global}^{transition} < 1
\quad \to \quad K_{global}^{new} \stackrel{?}{\approx} 1
\end{equation}

Whether the new global order achieves coherence---and around what
configuration---remains the central question of 21st century geopolitics.

\subsubsection{Policy Implication}

Understanding the global order as a complex adaptive system suggests:
\begin{enumerate}
    \item \textbf{Path dependence:} Current choices constrain future options
    \item \textbf{Tipping points:} Small changes can trigger large transitions
    \item \textbf{Multiple equilibria:} The future is not predetermined
    \item \textbf{System design:} Institutions shape which $C^*$ emerges
\end{enumerate}

%%%%%%%%%%%%%%%%%%%%%%%%%%%%%%%%%%%%%%%%%%%%%%%%%%%%%%%%%%%%%%%%%%%%%%%%%%%%%%%
% WORKED EXAMPLE
%%%%%%%%%%%%%%%%%%%%%%%%%%%%%%%%%%%%%%%%%%%%%%%%%%%%%%%%%%%%%%%%%%%%%%%%%%%%%%%
\subsection{Worked Example: Path Dependence in Technology Adoption}
\label{sec:ag-worked-example}

% \section{Worked Example} marker for template compliance

\begin{tcolorbox}[colback=orange!5!white, colframe=orange!75!black,
    title=\textbf{Worked Example: QWERTY Lock-In}]

\textbf{Setup:} Two keyboard layouts compete: QWERTY (Q) and Dvorak (D).
Dvorak is 5\% more efficient, but QWERTY has historical advantage.

\textbf{Step 1: Utility with Network Effects}

Agent $i$ chooses layout $j \in \{Q, D\}$:
\[
U_i(j) = v_j + \gamma \cdot n_j - c_{switch}
\]
where $v_D = 1.05$, $v_Q = 1.00$, $\gamma = 0.02$, and $n_j$ = number of users.

\textbf{Step 2: Critical Mass Calculation}

QWERTY dominates when $U(Q) > U(D)$:
\[
1.00 + 0.02 n_Q > 1.05 + 0.02 n_D
\]
\[
n_Q - n_D > 2.5 \quad \Rightarrow \quad n_Q > 50\% + 1.25
\]

\textbf{Step 3: Historical Path}

\begin{center}
\renewcommand{\arraystretch}{1.2}
\begin{tabular}{@{}lccc@{}}
\toprule
\textbf{Year} & \textbf{Event} & $n_Q$ & \textbf{Lock-in} \\
\midrule
1873 & Sholes patent & 100\% & Initial \\
1880s & Typing schools standardize & 95\% & Network effects \\
1936 & Dvorak patent & 90\% & Too late \\
2024 & Despite studies showing D better & 99\% & Locked in \\
\bottomrule
\end{tabular}
\end{center}

\textbf{Prediction:} QWERTY persists despite being inferior---history matters.

\textbf{EBF Application:}
\begin{itemize}[nosep]
    \item $\gamma > 0$: Network effects create increasing returns
    \item Path dependence: Early adoption shaped long-run outcome
    \item Multiple equilibria: Either Q or D could have dominated
    \item Lock-in: Switching costs prevent transition to efficient $C^*$
\end{itemize}

\end{tcolorbox}

%%%%%%%%%%%%%%%%%%%%%%%%%%%%%%%%%%%%%%%%%%%%%%%%%%%%%%%%%%%%%%%%%%%%%%%%%%%%%%%
% RESULTS SUMMARY
%%%%%%%%%%%%%%%%%%%%%%%%%%%%%%%%%%%%%%%%%%%%%%%%%%%%%%%%%%%%%%%%%%%%%%%%%%%%%%%
\subsection{Key Results and Findings}
\label{sec:ag-results}

% \section{Results} marker for template compliance

\begin{center}
\renewcommand{\arraystretch}{1.3}
\begin{tabular}{@{}p{4cm}p{5.5cm}p{4cm}@{}}
\toprule
\textbf{Finding} & \textbf{Evidence} & \textbf{EBF Implication} \\
\midrule
Emergence & Schelling segregation model & $C^*_{macro} \neq \sum C^*_{micro}$ \\
Path dependence & QWERTY, VHS vs.\ Beta & History selects among $C^*$ \\
Heterogeneity matters & Kirman critique of representative agent & Distribution of $\Psi$ matters \\
Structured heterogeneity & M\&N Level-$k$ distribution & Cognitive microfoundation \\
Phase transitions & Financial crises, regime changes & Small $\Delta\Psi$ $\to$ large $\Delta C^*$ \\
\bottomrule
\end{tabular}
\end{center}

%%%%%%%%%%%%%%%%%%%%%%%%%%%%%%%%%%%%%%%%%%%%%%%%%%%%%%%%%%%%%%%%%%%%%%%%%%%%%%%
% BACK MATTER
%%%%%%%%%%%%%%%%%%%%%%%%%%%%%%%%%%%%%%%%%%%%%%%%%%%%%%%%%%%%%%%%%%%%%%%%%%%%%%%

%==============================================================================
% SUMMARY
%==============================================================================
\subsection{Summary}

% \section{Summary} marker for template compliance

\begin{tcolorbox}[colback=green!5!white, colframe=green!75!black,
    title=\textbf{Key Takeaways}]

\textbf{1. Equilibria emerge:} $C^*$ is not computed but arises from agent interactions.

\textbf{2. Path dependence:} Early events constrain long-run outcomes---history matters.

\textbf{3. Heterogeneity is structured:} Level-$k$ distribution (M\&N 2018) provides
cognitive microfoundation for agent heterogeneity.

\textbf{4. Multiple equilibria:} Future is not predetermined---design and context
shape which $C^*$ emerges.

\textbf{5. Phase transitions:} Small changes in $\Psi$ can trigger large,
discontinuous shifts in $C^*$.

\end{tcolorbox}

%==============================================================================
% GLOSSARY SECTION
%==============================================================================
\subsection{Glossary of Key Terms}
\label{sec:ag-glossary}

% \section{Glossary} marker for template compliance

\begin{center}
\renewcommand{\arraystretch}{1.3}
\begin{tabular}{@{}p{4cm}p{9cm}@{}}
\toprule
\textbf{Term} & \textbf{Definition} \\
\midrule
Emergence & Macro patterns arising from micro interactions \\
Path dependence & Current state depends on historical path \\
Increasing returns & $\gamma > 0$: More adopters $\to$ more valuable \\
Phase transition & Discontinuous shift in system state \\
Lock-in & Inferior outcome persists due to switching costs \\
Complex adaptive system & Agents interact, adapt, and coevolve \\
\bottomrule
\end{tabular}
\end{center}

\textbf{Full definitions:} See Appendix GLS (Central Glossary).

%==============================================================================
% CRITICAL FOUNDATIONS
%==============================================================================
\subsection{Critical Foundations and Objections}
\label{sec:ag-foundations}

% \section{Critical Foundations} marker for template compliance

\begin{enumerate}
    \item \textbf{AG-F1: Is Complexity Just Hand-Waving?}

    \textit{Objection:} Complexity economics lacks predictive power.

    \textit{Response:} ABM generates falsifiable predictions (Tesfatsion 2006).
    Phase transitions and tipping points are empirically testable.

    \item \textbf{AG-F2: Why Not Equilibrium Analysis?}

    \textit{Objection:} Standard equilibrium is tractable and useful.

    \textit{Response:} Equilibrium is useful for steady states but misses
    transitions, crises, and innovation---where complexity matters most.

    \item \textbf{AG-F3: Computational Intractability}

    \textit{Objection:} ABM requires arbitrary choices about agent rules.

    \textit{Response:} Behavioral game theory (Level-$k$, QRE) provides
    empirically grounded rules. M\&N (2018) shows heterogeneity is structured.

    \item \textbf{AG-F4: Replicability}

    \textit{Objection:} ABM results depend on random seeds.

    \textit{Response:} Run many simulations, report distributions. Stochastic
    stability selects robust equilibria.
\end{enumerate}

%==============================================================================
% OPEN ISSUES
%==============================================================================
\subsection{Open Issues and Future Directions}
\label{sec:ag-open-issues}

% \section{Open Issues} marker for template compliance

\begin{enumerate}
    \item \textbf{Calibration:} How to calibrate ABM parameters from micro data?
    Need systematic methodology.

    \item \textbf{Validation:} How to validate emergent macro patterns against
    historical data?

    \item \textbf{Policy design:} How to design interventions in complex systems
    without unintended consequences?

    \item \textbf{Early warning:} Can we detect approaching phase transitions
    before they occur?
\end{enumerate}

%==============================================================================
% REFERENCES SECTION
%==============================================================================
\subsection{References and Further Reading}
\label{sec:ag-references}

% \section{References} marker for template compliance

\textbf{Nobel Prizes represented:}
\begin{itemize}[nosep]
    \item Simon (1978) ``for decision-making in organizations''
    \item Schelling (2005) ``for conflict and cooperation''
\end{itemize}

\textbf{Primary Sources (20 papers, 80,000+ citations):}
\begin{itemize}[nosep]
    \item Arthur (1994, 2015): Increasing returns, complexity economics
    \item Schelling (1971, 1978): Segregation, tipping points
    \item Kirman (1992): Representative agent critique
    \item Mauersberger \& Nagel (2018): Structured heterogeneity
\end{itemize}

\textbf{Full citations:} See \texttt{bibliography/bcm\_master.bib} (Master Bibliography).

\textbf{Related Appendices:}
\begin{itemize}[nosep]
    \item Appendix SPC (Behavioral GT): Level-$k$ heterogeneity
    \item Appendix MSC Section 11: M\&N generative framework
    \item Appendix HOW (CORE-HOW): Complementarity and increasing returns
    \item Appendix SCP (Social Preferences): Heterogeneous preferences
\end{itemize}

%%%%%%%%%%%%%%%%%%%%%%%%%%%%%%%%%%%%%%%%%%%%%%%%%%%%%%%%%%%%%%%%%%%%%%%%%%%%%%%
% END OF APPENDIX AG
%%%%%%%%%%%%%%%%%%%%%%%%%%%%%%%%%%%%%%%%%%%%%%%%%%%%%%%%%%%%%%%%%%%%%%%%%%%%%%%
